\section*{A row-marginal-floor closure theorem for three address bits}

All logarithms are base two.  Let
\[
 f:\{-1,1\}^{A}\times\{-1,1\}^{Y}\to\{-1,1\},\qquad |A|=3,
\]
and write
\[
 p(T,U)=\widehat f(T,U)^2,\qquad
 w_T=\sum_U p(T,U),\qquad q_T(U)=p(T,U)/w_T
\]
when $w_T>0$.  Put
\[
 \theta_{T,j}=\Pr_{U\sim q_T}[j\in U].
\]
Thus
\[
 I_A(f)=\sum_T|T|w_T,
 \qquad I_Y(f)=\sum_Tw_T\sum_j\theta_{T,j}.             \tag{1}
\]

We use the already proved complete three-address marginal ledger
\[
 H(w)\le 4I_A(f)+8I_Y(f).                               \tag{2}
\]
It follows from the exact three-bit Boolean-fiber trichotomy and a single
truth-table label payment.  No row-product hypothesis enters (2).

\paragraph{Theorem 1 (marginal-floor closure).}
Assume that for some $0<\beta\le1$, every $T$ with $w_T>0$ and every
$j\in Y$ satisfies
\[
 \theta_{T,j}=0\quad\hbox{or}\quad \theta_{T,j}\ge\beta. \tag{3}
\]
Then
\[
 \boxed{
 H(f)\le4I_A(f)+\left(8+{h_2(\beta)\over\beta}\right)I_Y(f).} \tag{4}
\]
In particular this is dimension independent for every fixed $\beta$.
Neither independence of the coordinates under $q_T$, nor a common row law,
nor a sign/phase hypothesis is required.

\paragraph{Proof.}
The full squared-spectrum chain rule is
\[
 H(f)=H(w)+\mathcal R_Y(f),\qquad
 \mathcal R_Y(f)=\sum_Tw_TH(q_T).                        \tag{5}
\]
Entropy subadditivity for the indicator vector of $U\sim q_T$ gives
\[
 H(q_T)\le\sum_j h_2(\theta_{T,j}).                      \tag{6}
\]
The function $h_2(x)/x$ is strictly decreasing on $(0,1)$, since
\[
 {d\over dx}{h_2(x)\over x}={\log_2(1-x)\over x^2}<0.  \tag{7}
\]
Consequently (3) implies
\[
 H(q_T)\le {h_2(\beta)\over\beta}\sum_j\theta_{T,j}.
\]
After multiplying by $w_T$, summing in $T$, and using (1),
\[
 \mathcal R_Y(f)\le {h_2(\beta)\over\beta}I_Y(f).       \tag{8}
\]
Equations (2), (5), and (8) prove (4). \qed

\paragraph{Corollary 2 (fair/frozen row closure).}
If every row marginal belongs to $\{0,1/2,1\}$, then
\[
 \boxed{H(f)\le4I_A(f)+10I_Y(f)\le10I(f).}             \tag{9}
\]
This covers correlated row laws: for example, a row law uniform on an
arbitrary affine spectral subspace has only deterministic or fair coordinate
marginals, whether or not that subspace is a coordinate cell.  Arbitrary
Fourier phases are allowed.  Thus (9) is strictly a marginal theorem, not a
hidden product or orthogonality assumption.

The prior row-monomial theorem had $\mathcal R_Y=0$ and constant eight.
Corollary 2 enlarges it to rows with exponentially many, mutually correlated
Fourier atoms, at the honest combined constant ten.  The coefficient ten
does not pretend that the $8I_Y$ truth-table-label charge and the $2I_Y$ row
entropy charge are disjoint; both terms are retained in the final budget.

\paragraph{A growing non-monomial family.}
Let $h$ be any three-bit Boolean gate.  On mutually disjoint even-dimensional
blocks $Y_i$, let $b_i$ be a bent Boolean function, and set
\[
 F(a,y)=h(a_1b_1(y_1),a_2b_2(y_2),a_3b_3(y_3)).         \tag{10}
\]
Expanding $h$ shows that row $T$ is
$\widehat h(T)\prod_{i\in T}b_i$.  Its squared spectrum is fair on every
coordinate in $\bigcup_{i\in T}Y_i$ and frozen elsewhere, so (9) applies in
arbitrary data dimension.  In fact the exact formulas are
\[
 H(F)=H(h)+\sum_i |Y_i|\,\operatorname{Inf}_i(h),
 \quad
 I(F)=I(h)+\frac12\sum_i |Y_i|\,\operatorname{Inf}_i(h). \tag{11}
\]
The three-bit fiber inequality $H(h)\le4I(h)$ therefore gives the sharper
$H(F)\le4I(F)$ for this explicit family.

\paragraph{Why the parameter cannot simply be deleted.}
The scalar cost in (8) satisfies
\[
 {h_2(\beta)\over\beta}=\log_2(1/\beta)+\log_2 e+o(1)
 \qquad(\beta\downarrow0).                              \tag{12}
\]
This divergence is sharp for entropy subadditivity, already for a one-bit
Bernoulli row law.  Actual Boolean compatibility can sometimes pay such a
rare row through address energy or arithmetic granularity, but that payment
is not contained in (6).  Indeed, on one data bit a genuine four-bit point
defect produces a three-address row marginal $1/50$, so the fair/frozen
hypothesis is not automatic even in the smallest extension.

The next unpaid lemma is therefore a Boolean-coupled rare-marginal payment:
replace the minimum $\beta$ in (4) by address/granularity reserves without
reusing the $8I_Y$ complete-label payment.  A fixed-$K_3$ result would still
leave the general-$k$ truth-table distance loss $2^k$ unresolved.

\paragraph{Scope.}
Theorem 1 and Corollary 2 are proper-class FEI theorems.  They do not prove
the unrestricted conditional-product $K_3$ theorem, a bounded recurrence in
$k$, or full FEI.
